\documentclass[11pt]{article}
\usepackage[margin=1in]{geometry}
\usepackage[utf8]{inputenc}
\usepackage[T1]{fontenc}
\usepackage{hyperref}
\hypersetup{colorlinks=true, linkcolor=blue, citecolor=blue, urlcolor=blue,
  pdftitle={The Late Channel}, pdfauthor={Caio Vicentino}}
\emergencystretch=3em \sloppy \hyphenpenalty=2000 \exhyphenpenalty=0 \hbadness=10000
\usepackage{booktabs}\usepackage{amsfonts}\usepackage{amsmath}\usepackage{amssymb}
\usepackage{xcolor}\usepackage{array}\usepackage{graphicx}

\title{The Late Channel: A Single Late Layer Band Carries the Causal\\
Content of Action \emph{and} Reasoning in a 27B Agent\\[2pt]
\large Mechanistic chain-of-thought faithfulness as a monitoring signal}
\author{%
  Caio Vicentino \\ OpenInterpretability \\ Fortaleza, Brazil \\
  ORCID: 0009-0003-4331-6259 \\ \texttt{caio@openinterp.org}}
\date{June 2026}

\begin{document}
\maketitle

\begin{abstract}
Chain-of-thought (CoT) monitorability is a leading safety bet for reasoning agents, but the question is
\emph{behavioral}: does the written CoT predict behavior? We ask the \emph{mechanistic} version on an
open-weight 27B reasoning model (Qwen3.6-27B) with a full-stack sparse autoencoder (SAE; 11 layers,
$d_{\text{sae}}{=}40960$): are the SAE features active at the post-reasoning decision point \emph{causal}
for what the model does, and \emph{where} do they live? Using a single causal-patch test
(decode--re-encode of the CoT decision-state into a no-CoT run, controlled for SAE reconstruction error and
against a random-feature baseline), we find a consistent picture across three decisions. (i) For the
final \emph{answer} to hard multi-step problems ($n{=}60$, 16 CoT-flips), patching the CoT features recovers
the reasoned answer with $\Delta\log p = +2.72\,[+2.18,+3.28]$ at L59, with the effect $\approx 0$ through
L47 (\textsc{late} $+2.13$ vs \textsc{early} $-0.02$, disjoint 95\% CIs; 16/16 consistent). (ii) For an
\emph{agent action} (tool call) in trap scenarios ($n{=}32$, 10 flips), the same test gives \textsc{late}
$+1.65\,[+1.17,+2.12]$ vs \textsc{early} $+0.08$, peaking at L63. (iii) The same late band (L51--63) is where
a prior result localizes \emph{action commitment} itself. Faithfulness is \textbf{conditional}: causal when
the CoT changes the outcome, $\approx 0$ (performative) when the model already knew (44/60 correct with no
CoT). We argue the late band is the natural locus for mechanistic action/answer monitoring---answering
DeepMind's call to ``inspect the model's inner workings''---while noting (from a companion result) that this
locus is \emph{not} adversarially robust as a control point.
\end{abstract}

\section{Introduction}
Reasoning models expose a chain of thought, and a growing safety program proposes to \emph{monitor} it: read
the CoT, flag bad behavior \cite{cotmon2025,baker2025}. But the CoT is text, and the field's own results show it
is an unreliable readout of the underlying computation: models can be unfaithful \cite{turpin2023,anthropic2025},
and monitorability degrades under optimization pressure \cite{cotmon2025,baker2025}. The position paper that
defines the agenda explicitly defers the mechanistic route as ``requiring further breakthroughs''
\cite{cotmon2025}. Concurrently, Google DeepMind's agent-security roadmap states that ``reading their verbalized
reasoning will not be enough'' and that securing agents ``will need to \ldots inspect the model's inner
workings'' \cite{deepmind2026}.

This paper takes the mechanistic route. We ask, on a single open-weight 27B reasoning model with a full-stack
SAE: \emph{are the features active during/after reasoning causal for the model's answer and action, and where
do they live?} Our contribution is a unifying empirical law---\textbf{the late channel}---and a method.

\paragraph{Contributions.}
\begin{itemize}
  \item A single causal-patch test (\S\ref{sec:method}) for mechanistic faithfulness that controls for SAE
  reconstruction error and a random-feature baseline---an application of a causal-rigor protocol from prior work.
  \item \textbf{Reasoning faithfulness is causal and late} (\S\ref{sec:tierA}): CoT decision-state features
  causally drive the answer (peak $+2.72\log p$ at L59), $\approx 0$ before L51, and \emph{conditional}---causal
  only when the CoT changes the answer; performative otherwise.
  \item \textbf{It holds for agent actions} (\S\ref{sec:tierB}): the same late, causal pattern governs tool-call
  decisions in agent-trap scenarios.
  \item \textbf{The three-faced law} (\S\ref{sec:unified}): the same late band (L51--63) carries the causal
  content of action commitment, reasoned answers, and reasoned agent actions---a single monitoring locus.
  \item An honest boundary (\S\ref{sec:disc}): the late band is a good \emph{monitor} but, per a companion
  result, \emph{not} an adversarially-robust \emph{control} point.
\end{itemize}

\section{Related work}
\textbf{Behavioral CoT faithfulness} is well studied and behavioral-only: unfaithfulness via input-perturbation
\cite{turpin2023}, reasoning-model hint-omission \cite{anthropic2025}, in-the-wild measurement, and the
monitorability position \cite{cotmon2025}---none open the model. \textbf{Mechanistic faithfulness} via
SAEs+activation patching has been shown on \emph{Pythia-2.8B base} \cite{chen2025}, finding causality is
scale-dependent (null at 70M, causal at 2.8B); we extend this to a 27B native \emph{reasoning} model.
\textbf{Internal monitoring beats text}: residual-stream probes match/beat a full-trace CoT monitor
\cite{mirtaheri2026}, and attention probes predict the answer early (``performative reasoning'')
\cite{goodfire2026}---both \emph{correlational}; we test \emph{causally} with named SAE features and on agent
\emph{actions}. \textbf{The action lever} (this arc): control of the \texttt{finish}/commit decision lives in a
late action-commitment band, not the mid-layer verdict \cite{leverlate,levergen}.

\section{Method}\label{sec:method}
\textbf{Model \& SAE.} Qwen3.6-27B (64 layers, reasoning/``thinking'' mode), bf16. A full-stack TopK SAE
($d_{\text{sae}}{=}40960$, $k{=}128$) trained on 11 layers $\{15,19,23,27,35,39,43,47,51,59,63\}$; encode
$z=\mathrm{TopK}(\mathrm{ReLU}((x-b_{\text{dec}})W_{\text{enc}}+b_{\text{enc}}))$, decode $\hat x=zW_{\text{dec}}+b_{\text{dec}}$.

\textbf{The causal-patch test.} For a decision, we run the model \emph{with} CoT (it thinks, then a forced slot
elicits the answer/tool token) and \emph{without} CoT (direct answer). We keep items where the CoT
\emph{flips} the no-CoT outcome (so the reasoning does work). At each SAE layer $L$, we encode the residual at
the CoT decision point to $z_{\text{cot}}$, and in the \emph{no-CoT} forward replace the decision-point residual
with $\hat x(z_{\text{cot}})$; we measure $\Delta\log p$ of the CoT-outcome token. The clean causal signal is
$d_{\text{cot}}-d_{\text{self}}$, where $d_{\text{self}}$ patches the no-CoT features' own reconstruction
(controls for SAE recon error); the noise control is $d_{\text{rnd}}-d_{\text{self}}$ with a random
top-$k$ feature set of matched magnitude. Outcomes are confirmed behaviorally (forced-slot argmax), paired per
item, with bootstrap 95\% CIs over flipped items. This is the causal-rigor protocol of \cite{leverlate} applied
to features.

\section{Reasoning faithfulness is causal and late}\label{sec:tierA}
We use 60 hard multi-step problems (serial arithmetic / counterintuitive / code traps) as MCQs with options
shuffled per item; 16 flip the no-CoT answer; in 44/60 the model is \emph{already correct without any CoT}.
Table~\ref{tab:tierA} reports the per-layer causal signal over the 16 flipped items.

\begin{table}[h]\centering\small
\begin{tabular}{lrr}
\toprule
Layer & causal $d_{\text{cot}}-d_{\text{self}}$ [95\% CI] & noise $d_{\text{rnd}}-d_{\text{self}}$\\
\midrule
L15--L43 & $\approx 0$ (CIs cross 0) & $-8$ to $-10$\\
L47 & $+0.17\,[+0.03,+0.29]$ & $-10.5$\\
L51 & $+1.55\,[+1.28,+1.85]$ & $-11.3$\\
\textbf{L59} & $\mathbf{+2.72\,[+2.18,+3.28]}$ & $-10.5$\\
L63 & $+2.12\,[+1.62,+2.61]$ & $-12.1$\\
\midrule
\textsc{late} (L51--63) & $+2.13\,[+1.83,+2.41]$ & \\
\textsc{early} (L15--27) & $-0.02\,[-0.05,+0.01]$ & \\
\bottomrule
\end{tabular}
\caption{Tier-A (answers, $n{=}60$, 16 flips). The causal signal is $\approx 0$ through L47 and concentrates in
L51--63 (\textsc{late} vs \textsc{early} CIs disjoint); L59 causal $>0$ in 16/16 flipped items (min $+0.60$, max
$+4.81$). Random patches destroy the answer at every layer.}\label{tab:tierA}
\end{table}

\textbf{Faithfulness is conditional.} The patch only matters where the CoT changes the answer. In the 44/60
items the model already answers correctly without thinking, the late causal signal is small---the written CoT is
\emph{performative}, the mechanistic version of \cite{goodfire2026}'s ``reasoning theater'', now measured causally.

\section{It holds for agent actions}\label{sec:tierB}
We repeat the test where the ``answer'' is a committed tool call. 32 agent-trap scenarios across four domains
(crypto/files/db/deploy): the surface instruction suggests an irreversible action
(\texttt{send/delete/drop/deploy}); a detail suggests a safe one. 10 flip. Table~\ref{tab:tierB} shows the same
late, causal pattern.

\begin{table}[h]\centering\small
\begin{tabular}{lr}
\toprule
Layer & causal $d_{\text{cot}}-d_{\text{self}}$ (10 flips)\\
\midrule
L15--L47 & $\approx 0$ ($-0.14$ to $+0.27$)\\
L51 & $+0.84$ \\
L59 & $+1.89$ \\
L63 & $+2.21$ \\
\midrule
\textsc{late} (L51--63) & $+1.65\,[+1.17,+2.12]$\\
\textsc{early} (L15--27) & $+0.08\,[-0.14,+0.32]$\\
\bottomrule
\end{tabular}
\caption{Tier-B (agent actions, $n{=}32$, 10 flips). Same late-band concentration; CIs disjoint.}\label{tab:tierB}
\end{table}

\textbf{Behavioral side-finding.} With CoT the agent chose the \emph{cautious} (safe) action in 27/32 traps
(only 5 irreversible)---honest reasoning acts as a safeguard on these scenarios. We note a limitation: because
the irreversible class is small (5/32), the per-layer action-\emph{prediction} AUROC, although late-dominant and
above a random-feature baseline, rests on too few positives to be conclusive (future work).

\section{The three-faced law}\label{sec:unified}
The same late band (L51--63) carries the causal content of three distinct things on this model: (i) action
commitment itself \cite{leverlate,levergen}; (ii) the reasoned answer (\S\ref{sec:tierA}); (iii) the reasoned
agent action (\S\ref{sec:tierB}). Decision-relevant information consolidates in a narrow late window, $\approx 0$
in mid layers---a single, consistent locus for reading what the model is about to do.

\section{Discussion}\label{sec:disc}
\textbf{A monitoring locus, not a control point.} The late channel is where mechanistic monitoring should read:
it carries the causal content and (per \cite{mirtaheri2026,goodfire2026}) predicts behavior earlier and more
reliably than CoT text. This directly instantiates DeepMind's ``inspect inner workings'' \cite{deepmind2026} for
a reasoning agent. \emph{However}, monitoring is not control: a companion experiment shows the same late
action-commitment band, used as a brake, collapses under an adaptive white-box adversary (attack success
$0\to1.0$ at a small budget that knows the locus, while a norm-matched random perturbation does nothing). The
honest reading: the late channel is a strong \emph{auditing} signal under a non-adversarial threat model, not an
adversarially-robust safeguard.

\section{Limitations}\label{sec:lim}
Curated stimulus sets (serial-computation MCQs; hand-built agent traps), not standard benchmarks; modest flip
counts ($16$ and $10$)---though effects are large and 16/16 consistent in Tier-A. Single model. The patch
substitutes the last-token decision residual. The prediction/monitoring AUROC (Tier-B) is underpowered
(5 positives). Causal results are conditional on the CoT changing the outcome; the performative regime is the
common case here.

\section{Conclusion}
On a 27B reasoning agent, a single late layer band carries the causal content of what the model answers and
what it does, and that content is faithful (causal) exactly when the reasoning changes the outcome. The late
channel is the natural locus for mechanistic monitoring of reasoning agents---and a sharper target than CoT
text---while remaining, as a control point, adversarially brittle.

\begin{thebibliography}{99}\small
\bibitem{turpin2023} Turpin et al. Language Models Don't Always Say What They Think. NeurIPS 2023. arXiv:2305.04388.
\bibitem{anthropic2025} Chen et al. (Anthropic). Reasoning Models Don't Always Say What They Think. 2025. arXiv:2505.05410.
\bibitem{cotmon2025} Korbak et al. Chain-of-Thought Monitorability. 2025. arXiv:2507.11473.
\bibitem{baker2025} Baker et al. (OpenAI). Monitoring Reasoning Models for Misbehavior\dots 2025. arXiv:2503.11926.
\bibitem{chen2025} Chen, Plaat, van Stein. How Does Chain-of-Thought Think? \dots SAE+patching. 2025. arXiv:2507.22928.
\bibitem{mirtaheri2026} Mirtaheri \& Belkin. Detecting motivated reasoning from internal representations. 2026. arXiv:2603.17199.
\bibitem{goodfire2026} Goodfire. Reasoning Theater\dots performative reasoning. 2026. arXiv:2603.05488.
\bibitem{deepmind2026} Google DeepMind. Securing the Future of AI Agents. 2026.
\bibitem{leverlate} Vicentino. The Lever Is Late. 2026. Zenodo 10.5281/zenodo.20534219.
\bibitem{levergen} Vicentino. The Lever Generalizes---and It Brakes. 2026. Zenodo 10.5281/zenodo.20634838.
\end{thebibliography}
\end{document}
